\section*{Abstract}
\addcontentsline{toc}{section}{Abstract}

Roadworks and scheduled maintenance force city authorities to close road
segments and publish deviation plans, yet the choice between candidate detours
is usually made by expert judgement and inspected only qualitatively. This
thesis develops a reproducible experimental framework that turns that choice
into a measurable comparison. From a SUMO road network and a Traffic Analysis
Zone layer, the framework synthesises network-aware demand using a gravity
origin-destination model, generates candidate detours around a closed edge
(manually, with Yen's $k$-shortest loopless paths, and with an upstream search
that intercepts traffic before it reaches the closure), rewrites the routes of
affected vehicles, and evaluates every candidate against a paired baseline over
repeated Monte Carlo simulations with randomised route assignment. A decision
layer rejects infeasible plans, screens the rest for Pareto dominance and ranks
the survivors by additive weighted scoring, using Dirichlet weight sampling to
test how far the ranking depends on the chosen weights.

The framework is exercised on three scenarios of increasing realism: a synthetic
\texttt{city\_tiny} grid, an OpenStreetMap topology around Flagey in Brussels and a closure of the Delta tunnel in Auderghem under both night-time and all-day
regimes. On the synthetic network, automated generation widened the candidate
set beyond the manually defined corridors, surfacing both shorter replacement
paths and interception points upstream of the closure itself. On both
real topologies the plan minimising aggregate delay was not the one distributing
the burden most evenly, so the recommendation held in only about three quarters
of sampled weight vectors (a trade-off the decision layer exposes rather than
hides behind a single number). Re-scoring the Delta candidates under full-day
exposure left the recommended plan unchanged but widened the worst-to-best delay
ratio from $2.6$ to $9.6$, showing that candidates must be evaluated at the
demand they will actually experience. All demand is synthetic, so the absolute
figures are not calibrated predictions for any real city.
